\chapter{IoT Systems}

An IoT system consists of all the HW and SW components that are required to implement an IoT application. In order to reduce the complexity of an IoT system, developpers tend to reuse the same components across different applications, and here is where the concept of \emph{IoT platforms} comes into play.

\section{IoT Platforms}

An IoT platform is a framework that provides a set of tools and services to facilitate the development, deployment, and management of IoT applications. By using an IoT platform, developers can focus on the application logic rather than dealing with the complexities of the underlying infrastructure. There are many IoT platforms available in the market, each with its own strengths and weaknesses. There are two main categories of IoT platforms:
\begin{itemize}
    \item \ul{Horizontal}: These are general-purpose platforms that can be used for a wide range of IoT applications. They provide a broad set of features and services that can be customized to fit the specific needs of different applications.
    \item \ul{Vertical}: These platforms are designed for specific industries or use cases. They provide specialized features and services that are tailored to the requirements of a particular domain.
\end{itemize}

\section{IoT System Architecture}

An IoT system architecture defines whole IoT system. Regarding the HW, it consists of 4 different types of components:
\begin{enumerate}
    \item \ul{IoT Things}: Also known as Smart Objects, these are the physical devices that actually interact with the environment.
    \item \ul{Gateways}: These are devices that act as intermediaries between the IoT Things and the Internet. Given that the IoT Things are often resource-constrained and may not have the capability to connect directly to the Internet, IoT Gateways provide the necessary connectivity and protocol translation.
    \item \ul{Internet}: This is the communication infrastructure that allows the IoT Gateways to connect and exchange data with other devices and services.
    \item \ul{Online Services}: These are cloud-based platforms and services that provide different functionalities to support the IoT applications.
\end{enumerate}

However, SW components of an IoT system are also crucial to the overall functionality and performance of the system. There are three main aspects from the SW perspective that affect the whole IoT system.

\subsection{SW Functional Blocks}

From the SW perspective, it is essential to identify the functional blocks that are required to implement an IoT application. It is difficult to define a standard set of functional blocks that can be applied to all IoT applications, but for the aim of this course, we will consider three functional blocks:
\begin{enumerate}
    \item \ul{IoT Networking}: SW necessary to enable communication between all the components of the IoT system.
    \item \ul{IoT Data Management}: SW responsible for measuring, collecting, storing, and making data available for querying and analysis.
    \item \ul{IoT Data Processing}: SW responsible for analyzing the data collected from the IoT Things and extracting useful insights and information.\\
    
    There are two main recipient groups of the data in order to make use of it:
    \begin{itemize}
        \item Humans: The data is presented to the users in a human-readable format, such as dashboards, reports, or visualizations.
        \item Algorithms: The data is used as input for machine learning models, predictive analytics, or other algorithms that can make decisions or take actions based on the data.
    \end{itemize}

    There are four main levels of data processing that can be applied to the data collected from the IoT Things:
    \begin{enumerate}
        \item \ul{Descriptive}: The data is only used to observe and describe the current state of the system.
        \item \ul{Diagnostic}: The data is used to identify the causes of certain events or conditions in the system and explain why they happened.
        \item \ul{Predictive}: The data is used to anticipate future events or conditions in the system based on historical data and patterns.
        \item \ul{Prescriptive}: The data is used to perform actions or make decisions based on the insights derived from the data analysis. This has several legal and ethical implications, and usually requires human supervision and intervention to ensure that the actions taken are appropriate and that someone is legally responsible for the consequences of those actions.
    \end{enumerate}
\end{enumerate}

\subsection{SW Blocks Placement}

The SW functional blocks can be executed in different specific HW components within the IoT system, depending on the specific requirements and constraints of the application. There are three main approaches for the placement of the SW functional blocks.
\begin{enumerate}
    \item \ul{Cloud-based IoT Systems}: The ``classic'' approach.
    
    In this approach, the biggest part of the three SW functional blocks are executed in the cloud.
    \begin{itemize}
        \item IoT Things are responsible for collecting data and sending it to the gateway. Therefore, they only execute a small part of the IoT Networking and IoT Data Management blocks.
        \item IoT Gateways are responsible for receiving the data from the IoT Things and forwarding it to the cloud. Therefore, they execute a small part of the IoT Networking.
        \item The online services in the cloud are responsible for executing the majority of the three SW functional blocks, including IoT Networking, IoT Data Management, and IoT Data Processing.
    \end{itemize}

    \item \ul{Edge-based IoT Systems}: The ``modern'' approach.
    
    This approach is similar to the cloud-based approach, but with a significant shift of the SW functional blocks towards the gateways. Given that the online services have limited resources and are under high load, it is necessary to offload some of the processing tasks to the gateways in order to reduce latency and improve performance.

    \item \ul{Mesh-based IoT Systems}: The ``future'' approach.
    
    In this approach, the SW functional blocks are even for shifted towards the IoT Things, which are becoming more powerful and capable of executing more complex tasks. This approach just considers gateways and online services as fallback options for the IoT Things in case they are not able to execute certain tasks.
\end{enumerate}

\subsection{SW Blocks Communication}

The SW functional blocks need to communicate with each other in order to exchange data and coordinate their actions. This communication can be achieved through different connectivity options and protocols, depending on the specific requirements and constraints of the application. We will explore them in more detail in the rest of the course.